\begin{proof}
\begin{enumerate}[label=\arabic*.]
\item An effective divisor dominating the base would restrict to
a nonempty effective divisor on a general complex torus.
Pairing its Chern class with the square of a Kähler form
on that torus is strictly positive. This contradicts
$H^2(M;\mathbb Z)=0$. Thus every divisor is vertical.
The pole divisor of a meromorphic function is vertical;
its restriction to a general compact fibre is holomorphic
and therefore constant. The graph and properness of $f$
descend it to a meromorphic function on $\mathbb P^1$.
Conversely pullbacks from the base occur. Hence the
meromorphic function field is $\mathbb C(t)$ and $a(M)=1$.
Every automorphism consequently induces a Möbius map on
the base. The two distinct multiplicities and the cusp
identify all three values separately, so this map is the
identity.

\item On smooth fibres an automorphism is affine. Its integral
linear part is locally constant and commutes with $M_1,M_2$.
Solving the two exact matrix equations in
\path{code/part1/analytic_automorphisms.py} gives the rational
centralizer $\mathbb Q I$; its integral invertible elements
are $\pm I$. At the order-five filling an affine lift with
linear part $\epsilon I$ would have to conjugate the same
finite affine generator. For its translation $b$ this
requires
\[
(I-M_1)b=(1-\epsilon)v_1/5\pmod\Lambda.
\]
Applying the invariant integral covector $\ell$ makes the
left side zero. For $\epsilon=-1$ the right side is
$2/5\pmod\mathbb Z$, which is nonzero. Thus every
automorphism is a translation.

\item Let $\mathcal A$ be the sheaf of translations extending to
the specified cusp model and finite root charts. There is
an exact analytic sheaf sequence
\[
0\to j_*\Lambda\to E_0\to\mathcal A\to0
\]
on the root orbifold. Here $j$ omits the cusp and the stalk
of $j_*\Lambda$ there is $\ker N$.
To check the last surjectivity at the cusp, lift a translation
to the toroidal cover. It preserves the fan precisely when
its valuation is an integral translation of the periodic
fan. Subtract the corresponding period; the remaining
characters are nonvanishing holomorphic functions on a
disc and possess logarithms. These give a local section
of $E_0$. The kernel consists exactly of the invariant
integral periods. This also proves that the sequence
accounts for all extendable translations, including
possible component permutations.

\item The cohomology of $j_*\Lambda$ in degree one is calculated
from three positive meridians. For the finite points use
the integer kernels of $I+M_i+\cdots+M_i^{m_i-1}$;
for the cusp use the primitive image of $N$.
Writing these column bases as $L_1,L_2,L_0$, the cocycle
and coboundary matrices are
\[
D_1=(L_1\ M_1L_2\ M_1M_2L_0),\qquad
D_0=\begin{pmatrix}L_1^{-1}(M_1-I)\\
L_2^{-1}(M_2-I)\\L_0^{-1}(M_0-I)\end{pmatrix}.
\]
Each inverse here denotes the unique coordinates in the
displayed integral column basis, not a square matrix
inverse. The exact $6$-by-$11$ and $11$-by-$6$ integer
matrices, their bases and Smith operations are exported
in \path{code/part1/analytic_automorphisms.json} and the
underlying parabolic program. They satisfy $D_1D_0=0$ and
give $\ker D_1/\im D_0=\mathbb Z/2$. The cocycle relations
are the finite-order relations and the product of the
three meridians; the cusp condition is precisely the
vanishing of its restriction modulo $(M_0-I)\Lambda$.
Thus this complex computes the required orbifold sheaf
cohomology, with its finite stabilizer torsion retained.

\item Now $H^0(j_*\Lambda)=0$ by P1.1,
$H^0(E_0)=\mathbb C$, and $H^1(E_0)=0$ by the floor
calculation. The exact sequence gives
\[
0\to\mathbb C\to H^0(\mathcal A)\to\mathbb Z/2\to0.
\]
The half-period
$\tau=(0,1,1,1,1,1)^t/2$ is fixed modulo $\Lambda$
by every monodromy matrix. It extends over the finite
quotients because the translations commute there.
At the cusp $N\tau$ is integral; primitiveness of $N$
provides $\lambda\in\Lambda$ with $N\lambda=N\tau$.
The representative $\tau-\lambda\in\ker N$ is a bounded
angular half-period and extends on every toric chart.
It is a nonzero order-two global section. Since the
additive group $\mathbb C$ has no nonzero torsion, it
splits the displayed exact sequence. Translations
commute, giving the asserted full group.

\item For an isomorphism between two members, the same intrinsic
base argument applies. Its integral linear part still
lies in the common centralizer. Either scalar $\pm I$
preserves a complex period plane exactly by equality
of the two planes. The graphs
$(\sigma+\xi\alpha e)(P)$ are equal precisely when their
parameters $\xi$ agree, since $e\notin F$ and $\alpha$
is nonzero. Translation does not change this condition.
This proves the parameter assertion.

\end{enumerate}
\end{proof}
